\section{The increment: new finite proof objects, not a second planner}

\subsection{What the repository already supplies}

At the audited revision, Forge is a design and research-prototype repository,
not an installed tactic. Its merged article already covers induction with
accumulator generalisation, polynomial recurrence and invariant synthesis,
several arithmetic witness engines, nonlinear certificates, demand-directed
Horn reasoning, proof-logging Boolean search, and the separation between search
and checked evidence. Its integration section already advocates a conservative
bridge through ordinary Lean goals and explicitly proved facts. Those ideas are
prerequisites here, not contributions of this article~\cite{forge-readme,forge-status,forge-structural,forge-integration}.

The useful next question is not ``which extra tactic should a scheduler try?''
It is: \emph{which large families of proof obligations admit a finite object
whose verification is much simpler than discovering the proof?} This question
leads to three workers with sharp input and output contracts.

\begin{table}[ht]
\centering\small
\begin{tabularx}{\linewidth}{@{}p{0.20\linewidth}YYY@{}}
\toprule
Worker & Input family & Returned object & New scope relative to the audited design\\
\midrule
Linear words & A finite-dimensional rational state model and an output identity & An invariant row-space simulation, or a separating word & All words, not just an additive polynomial sequence or one proposed invariant\\
Safe telescoping & A definite binomial-power sum & A recurrence plus a shifted-support flux & Hypergeometric summands and moving support, not polynomial antidifferencing\\
Recurrence transport & Two annihilating operators & A common left multiple and singularity/seed plan & Composition of nonconstant recurrences without silently cancelling singular leading terms\\
\bottomrule
\end{tabularx}
\caption{The three mathematical additions. A singularity plan still leaves explicit equality obligations.}
\end{table}

The workers also connect. A unary linear word model yields a constant-coefficient
recurrence. A telescoper yields a polynomial-coefficient recurrence. A common
left multiple gives both objects one recurrence, and singularity-aware induction
reduces equality to finitely many proved values. None of these connections
requires replacing \grind{}'s equality or arithmetic engines.

\subsection{A precise non-duplication claim}

The comparison was made against Forge's README, implementation-status ledger,
Leant/Djex transfer guide, structural-search and integration sections,
conclusion, and provenance appendix at commit
\texttt{674521027d968d59f7b83220ed52304a85cb55e2}. The provenance appendix maps the
merged proposals to their nine original submissions~\cite{forge-provenance}.
This is not a claim to have read every line of all nine frozen originals.
Repository code search returned no matches even for a term visibly present in
the sources, so negative search results were not used to establish absence.
The companion \code{docs/NON_DUPLICATION.md} records this limitation.

The underlying mathematics is not presented as new. Field-weighted automaton
equivalence and invariant-space algorithms are established methods; Kiefer's
notes give a particularly direct account~\cite{kiefer}. Creative telescoping
and Ore-operator computation likewise have mature implementations~\cite{holonomic,ore}.
The contribution is the concrete selection of fragments, the proof-object
interfaces, the boundary-aware implementation, and their proposed use inside
Forge. An existing polynomial-invariant engine could prove some overlapping
examples after a suitable invariant was supplied. What changes here is the
automatic reduction to a finite closure calculation, with a completeness theorem
for an explicitly stated fragment.

\subsection{What an actual gain would look like}

Consider a theorem about two different arithmetic folds over an arbitrary list.
A local saturation engine may simplify their definitions, but simplification
does not itself invent the relation connecting their hidden states. The linear
worker computes an invariant space in which that relation becomes a single
orthogonality check. Its proof is independent of list length.

Now consider a sum of cubes of binomial coefficients. A polynomial template in
$n$ is inappropriate: the sequence is not being claimed polynomial. The
summation worker instead proposes a recurrence and a flux whose finite
difference is the desired combination of summands. The sum disappears only
after both support boundaries have been justified.

Finally, two valid recurrence proofs need not use the same operator. Their
common left multiple provides a shared induction rule. If that operator has a
zero leading coefficient at a natural index, the proof must request an extra
value rather than pretend that its usual recurrence step still determines the
next term. This is a mathematical obligation, not a larger search budget.

\subsection{Evidence levels used in this article}

\status{All Python workers and tests described as executed were run. The Lean
source supplied in the bundle is an uncompiled semantic-lemma candidate. Lean,
Lake, and Elan were not available in the working runtime; attempted retrieval
routes did not establish an executable toolchain.}

A mathematical proof below establishes the stated theorem under its hypotheses.
A passing Python checker establishes that this implementation accepted a
particular encoded object. Neither alone establishes that an arbitrary Lean
expression was faithfully encoded, or that a Lean kernel accepted the resulting
proof. The proposed Lean acceptance path must discharge all of those separate
requirements. The raw ledger never upgrades ``plan accepted'' into ``theorem
proved''.

Forge's own conclusion correctly distinguishes additional prototype results
from closing its Lean-integration gap~\cite{forge-conclusion}. This article does
not erase that gap. Its first implementation milestone is correspondingly a
small matrix soundness theorem and source bridge, not three simultaneous large
integrations.
